\subsection{A graph presentation and an actual pullback}
The Fock algebra itself has a graph presentation, on a different graph from the original gentle quiver. Let $G_\alpha$ have vertices $q_h$ for $h\in H$ and sinks $b_e$ for $e\in Q_0$, and edges
\[
 s_h:q_h\to q_{\tau h},\qquad t_h:q_h\to b_{[\tau h]}.
\]
\begin{proposition}\label{thm:graphmodel}
$C^*(G_\alpha)\cong\mathcal T_A$, by $q_h\mapsto Q_{a_h}$, $b_e\mapsto v_e$, $s_h\mapsto L_{a_h}Q_{a_{\tau h}}$ and $t_h\mapsto L_{a_h}v_{[\tau h]}$. The integer graph relation at $h$ is $q_h-q_{\tau h}-b_{[\tau h]}$. Eliminating these relations recovers $K_1=\ker B_\Gamma$ and $K_0\cong\operatorname{coker}B_\Gamma\oplus\Z^{\#\operatorname{cyc}\sigma}$.
\end{proposition}
\begin{proof}
The two edges at $q_h$ have orthogonal initial pieces and orthogonal ranges summing to $Q_{a_h}$, by the Fock initial-projection identity. All vertex projections are nonzero. The length gauge acts with degree one on both edge families, so gauge-invariant uniqueness applies. They generate since $L_{a_h}=s_h+t_h$ and $P_e=v_e+\sum_{[h]=e}Q_{a_h}$. In the integer matrix a kernel coefficient is constant around each permitted cycle; its sink equations are exactly $B_\Gamma c=0$. In the cokernel, eliminate all but one state per cycle; its remaining relation is the sum of its incident sink generators. The remaining cycle generators are free, proving the stated groups.
\end{proof}
For the nodal interval, $G_\alpha$ has two loop vertices, each with an exit to the common vacuum sink. Chirvasitu--Hajac--Tobolski's actual pullback theorem \cite[Theorem3.4]{GraphPullbacks} gives
\[
 \{T\in\mathcal T_A:\pi(T)\in\C^2\}
 \cong C^*(E_1),\qquad \C^2\hookrightarrow C(\mathbb T)^2\text{ by constants},
\]
where $E_1$ has the same three vertices and countably many arrows from each loop vertex to the sink, with no loops. The arrow numbered $n\ge0$ maps to loop$^n$ followed by exit. This algebra is AF, with $K_0=\Z^3$, $K_1=0$. The extended edition verifies every hypothesis. The theorem does not turn arbitrary cycle-vertex identifications into pullbacks: its first graph is acyclic and its path functor must cover all pointed paths. Its sink-extension Theorem4.3 has further injection and path-coverage requirements.

\subsection{A separate diagonal-fiber C* completion retaining the pairing}
The formal order's residue gluing cannot be copied by taking diagonals of arbitrary full matrices: $\operatorname{diag}(E_{12}E_{21})\ne\operatorname{diag}(E_{12})\operatorname{diag}(E_{21})$. For $\ell>1$ there is no nonzero *-homomorphism from $M_\ell(C(\mathbb T))$ to a commutative C*-algebra, as restriction to its constant matrix units shows. Restrict the distinguished fiber first. Define
\[
 D_\ell=\{f\in C(\mathbb T,M_\ell(\C)):f(1)\text{ is diagonal}\},\qquad
 \mathcal D_H=\Bigl(\prod_cD_{|c|}\Bigr)\mathop{\times}_{\C^H}\C^{Q_0},
\]
where the first map evaluates diagonal entries at one and the second repeats $\lambda_e$ at the two darts paired by $\alpha$.
\begin{theorem}[A faithful polynomial pullback model]\label{thm:diagonalmodel}
There is a split exact sequence
\[
 0\to\bigoplus_c C_0(\mathbb T\setminus\{1\},M_{|c|})
 \to\mathcal D_H\to\C^{Q_0}\to0,
\]
so $K_0(\mathcal D_H)=\Z^{|Q_0|}$ and $K_1(\mathcal D_H)=\Z^{\#\operatorname{cyc}\sigma}$. The polynomial dessin algebra embeds faithfully and generates this C*-algebra with its adjoints. This embedding does not extend to a unital algebra map from all of $\Lambda$.
\end{theorem}
\begin{proof}
The section consists of the constant paired diagonal projections. The kernel is the displayed vanishing fiber, a sum of matrix suspensions, giving the groups by split exactness and suspension. For an explicit embedding set $g(z)=|1-z|^2/4$, choose one cut arrow per permitted cycle, and send $a_h$ to $g(z)z^{\epsilon_h}E_{h,\tau h}$ in its branch, where $\epsilon_h=1$ at the cut and zero otherwise. Send $e_e$ to the constant paired diagonal projection. Every forbidden product is zero. A full cycle at $h$ maps to $w_c(z)E_{hh}$, $w_c=g^{|c|}z$. This scalar function is injective on the circle: away from its unique zero at one its argument recovers $z$. Continuous functional calculus therefore supplies every $f(z)E_{hh}$ with $f(1)=0$. Multiplying by path images and approximating division away from one supplies every vanishing off-diagonal entry. Together with the constant paired diagonals these generate $\mathcal D_H$. Path normal forms in a fixed entry are a nonzero prefactor times a polynomial in $w_c$; its infinite range proves independence. Evaluation at one first separates the zero-length coefficients. Finally for $z_0\ne1$, $1-w_c(z_0)^{-1}c_h$ is a unit in the formal order, but its proposed image is singular at $z_0$. A unital map preserves units, proving nonextension.
\end{proof}
This is the C* algebra generated by a faithful polynomial representation and its adjoints, distinct from the completed order and from $\mathcal T_A$. The norm closure of the polynomial image alone need not contain those adjoints. Their $K$ groups already differ. The pairing belongs to the labelled presentation; an unmarked isomorphism class need not recover the ribbon ordering. No arithmetic descent, Frobenius action or completed-$L$ positive determinant is inferred from this additional valid gluing model.
